% title: 2026 春复变函数期末考试
% author: 贺丹青、王珺
% summary: 从香蕉空间旧站“试卷：复变函数”迁移的完整试卷题目。
% source: https://bananaspace.org/legacy/wiki/用户:Solution/_试卷:_复变函数
\section*{2026 春复变函数期末考试（贺丹青、王珺）}
\begin{enumerate}
    \item 判断(15分)\\
    \begin{itemize}
        \item[(1)]若函数$f(z)$在区域$\Omega$上全纯，则$f$必有原函数
        \item[(2)]设$f$在区域$\Omega$上全纯，若存在$\Omega$上互不相同的点列$\{z_k\}_n^{\infty}$使得$f(z_k)=0,k\in\mathbb{N}$，则$f(z)=0$
        \item[(3)]设$z=x+iy$，则$f(z)=e^x(3\sin y+3\sin x)$是全纯函数
        \item[(4)]$\{z:\operatorname{Re} z\geqslant0\}$不是$\mathbb{C}_{\infty}$中闭集
        \item[(5)]存在在$z=0$全纯的函数$f$使得对所有充分大的自然数$n$有$f\big(\frac1n\big)=\frac{1}{n^2}$
    \end{itemize}
    \item 计算(30分)
    \begin{itemize}
        \item[(1)]$\int_{|z|=1}\frac{1}{z^2(z^2-4)e^z}dz$.
        \item[(2)]计算$p(z)=z^4+6z^3+2z+2$在$\{z:|z|<1\}$内的零点个数.
        \item[(3)]求函数$f(z)=\frac{z^2-2z+5}{(z-2)(z^2+1)}$在圆环$\{1<|z|<2\}$上的Laurent展开式.
    \end{itemize}
    \item (15分)求双全纯函数$F:\mathbb{D}\cap\mathbb{H}\to L(0,\pi),L(0,\pi)=\{z\in\mathbb{C},0<\operatorname{Im} z<\pi\}$且$F(i(\sqrt2-1))=\frac{\pi i}{2}$，要求写出$F$的准确表达式，求说明这是由哪些初等共形映射复合而来.
    \item (15分)设$f$在$D^*(z_0,r)$上全纯，$z_0$是$f$的$n$阶极点，且是唯一极点，证明：存在$D^*(z_0,r)$上全纯函数$g$使得$z_0$是$g$的唯一的单极点，且满足$g^n(z)=f(z),z\in D^*(z_0,r)$.
    \item (10分)证明幂级数$\sum_{n=1}^{\infty}\frac{z^n}{n}$在$\{z\in\mathbb{C},|z|<1\}$上内闭一致收敛，但不一致收敛.
    \item (15分)设$f$在$\mathbb{D}=\mathbb{D}(0,1)$上解析，且满足$f(0)=f'(0)=f''(0)=0,|f''(z)|\leqslant 1,\forall z\in \mathbb{D}(0,1)$，证明$|f(z)|\leqslant\frac{|z|^3}{6},\forall z\in\mathbb{D}(0,1)$.
\end{enumerate}
